\chapter*{Abstract}\label{Chap:Abstract}
\addcontentsline{toc}{chapter}{Abstract}

This thesis addresses the position control of the ECP-730 magnetic levitation system in its attractive configuration, which is open-loop unstable, in the presence of the stick-slip effect produced by the contact between the magnet and its guide rod. The nonlinear model uses the manufacturer's electromagnetic force law and parameters identified experimentally in previous work; stick-slip is represented by Karnopp static friction and an actuator dead zone, simulated as a voltage hold. A linear controller with double integral action (PII) was designed on the linearized model, and its performance was evaluated by simulation in open loop, on the linear model with stick-slip, on the nonlinear model without and with this effect, in regulation and in tracking of four references (sinusoidal, trapezoidal, sawtooth and square-pulse), and in the stable configuration of the same device. The PII stabilizes the system in every test without explicitly compensating the friction or the dead zone, with practically zero mean error and, in the unstable configuration, without reaching the upper actuator limit except in brief transients after instantaneous reference changes. However, in simulation the levitated magnet does not come to rest on the reference: it oscillates in a limit cycle of 0.02 to 0.03~cm, caused by the interaction between integral action and static friction, a mechanism also confirmed through a first-principles derivation of the electromagnetic force. A comparison with a PI controller, a sweep over two families of PI and PII controllers, and a later joint sweep of the gain, low-frequency zero and lead network of both show that, under the tuning habitually used for both controllers, the second integral action reduces neither the limit-cycle amplitude nor the duration of the sticking phases; the amplitude depends mainly on the controller's gain, not on the number of integrators. That advantage of the second integration does appear, in bounded form, in a low-gain, larger-cycle regime outside the habitual tuning. It was also found that, with a maximum voltage of 3.5~V, the system can only change position in a regulated way above about $-3.6$~cm, which restricts the operating range.
